% ============================================================
\setpartpreamble{\partintro
  {Testen, versionieren, paketieren, und die Checkliste für den Weg von A bis Z.}
  {
    \item Kapitel~\ref{ch:test-strategy}: Teststrategie
    \item Kapitel~\ref{ch:versioning-and-compatibility}: Versionierung und Kompatibilität
    \item Kapitel~\ref{ch:packaging-and-documentation}: Packaging und Dokumentation
    \item Kapitel~\ref{ch:the-checklist-from-a-to}: Die Checkliste von A bis Z
  }}
\part{Qualität und Auslieferung}\label{part:quality}

% ------------------------------------------------------------
\chapter{Teststrategie}\label{ch:test-strategy}
\versionrisk{medium}

Plugins lassen sich erstaunlich weit ohne Datenbank testen. Modellbau, Type Mappings,
SQL-Übersetzung, Migrations-SQL und Snapshots brauchen nur einen Connection-String, keine
Verbindung. Die Teststrategie des Beispiels hat vier Ebenen:

\begin{datatable}{H{0.78}H{1.70}H{0.52}}
  \textbf{Ebene} & \textbf{Prüft} & \textbf{Datenbank} \\
  \tablehead
  Modelltests & Annotations, Conventions, Validierung, Optionen, Caching & nein \\
  \tablerule
  SQL-Tests & Übersetzungen per \code{To\allowbreak{}Query\allowbreak{}String},
  Migrations-SQL per Differ und Generator & nein \\
  \tablerule
  Snapshot-Tests & Code-Generator als String und als Roundtrip (Roslyn) & nein \\
  \tablerule
  Integrationstests & Echtes Verhalten: RLS, Rollen, Interceptors (Testcontainers) & ja
  (Docker) \\
\end{datatable}

\section{Testhilfen}

Tests bauen viele Contexts mit unterschiedlichem
\code{On\allowbreak{}Model\allowbreak{}Creating}. Weil EF das Modell pro internem Provider
cached, muss jeder Test einen eigenen Provider bekommen; dafür gibt es
\code{Enable\allowbreak{}Service\allowbreak{}Provider\allowbreak{}Caching(\allowbreak{}false)}:

\begin{sourcecode}[tests/\allowbreak{}.\allowbreak{}.\allowbreak{}.\allowbreak{}/\allowbreak{}Test\allowbreak{}Model.\allowbreak{}cs (\allowbreak{}gekürzt, vereinfacht)]{csharp}
public static TestContext Create(
    Action<TenancyOptionsBuilder>? tenancy = null,
    Action<ModelBuilder>? model = null,
    bool cacheServiceProvider = false,
    Action<DbContextOptionsBuilder>? configure = null)
{
    var builder = new DbContextOptionsBuilder<TestContext>()
        .UseNpgsql("Host=localhost;Database=unused")       // keine Verbindung nötig
        .UseAcmeTenancy(tenancy)
        .EnableServiceProviderCaching(cacheServiceProvider);
    configure?.Invoke(builder);
    return new TestContext(builder.Options, model);
}

public static IReadOnlyList<string> MigrationSql(DbContext? source, DbContext target)
{
    var operations = target.GetService<IMigrationsModelDiffer>().GetDifferences(
        source?.GetService<IDesignTimeModel>().Model.GetRelationalModel(),
        target.GetService<IDesignTimeModel>().Model.GetRelationalModel());

    return target.GetService<IMigrationsSqlGenerator>()
        .Generate(operations, target.GetService<IDesignTimeModel>().Model)
        .Select(c => c.CommandText.Trim()).ToList();
}
\end{sourcecode}

\begin{important}[Design-Time-Modell verwenden]
  Für Migrations-Tests immer das \textbf{Design-Time-Modell} verwenden
  (\code{IDesign\allowbreak{}Time\allowbreak{}Model}). \code{context.\allowbreak{}Model} ist
  das schlanke Laufzeitmodell; Annotations, die der Annotation-Provider nur zur Design-Zeit
  weitergibt, fehlen darin (Fehlerfall in
  Kapitel~\ref{ch:declarative-sql-from-annotations}).
\end{important}

\begin{note}[Warum nicht SQLite oder InMemory?]
  Tests gegen SQLite oder den InMemory-Provider sind verlockend, weil sie kein Docker
  brauchen. Für ein Plugin beweisen sie wenig: InMemory ist nicht relational, Type Mappings,
  Translators und Migrationen laufen also nie; SQLite hat einen anderen SQL-Dialekt, andere
  Type Mappings und keine Row Level Security. Ein provider-spezifisches Plugin wird mit
  seinem eigenen Provider getestet: ohne Datenbank (nur Connection-String) für alles bis zum
  SQL-Text, gegen die echte Datenbank für das Verhalten. Für Anwendungen, die das Plugin
  nutzen, sieht es anders aus (nächster Abschnitt).
\end{note}

\section{Testmodus für Anwendungen}

Anwendungen haben oft hunderte Tests ihrer Geschäftslogik auf SQLite oder InMemory. Sie prüfen
nicht die Isolation, sondern die eigene Logik, und sie dürfen nicht schon deshalb scheitern,
weil die Entitäten \code{ITenant\allowbreak{}Owned} implementieren. Das PostgreSQL-Paket
lehnt SQLite bewusst ab, und diese Prüfung abzuschalten hilft nicht: Das Paket ersetzt
Provider-Dienste durch Npgsql-Ableitungen, bildet \code{Tenant\allowbreak{}Id} auf
\code{uuid} ab und filtert mit \code{current\_\allowbreak{}setting(\allowbreak{})}. Der Kern
hat deshalb einen eigenen Testmodus, der nur die anwendungsseitige Hälfte des Plugins enthält:

% source: tests/Acme.EntityFrameworkCore.Tenancy.Tests/TestingModeTests.cs
\begin{sourcecode}[tests/\allowbreak{}.\allowbreak{}.\allowbreak{}.\allowbreak{}/\allowbreak{}Testing\allowbreak{}Mode\allowbreak{}Tests.\allowbreak{}cs (\allowbreak{}Auszug)]{csharp}
private sealed class AppContext : DbContext
{
    protected override void OnModelCreating(ModelBuilder modelBuilder)
        => modelBuilder.HasTenantQueryFilters(this);
}

new DbContextOptionsBuilder<AppContext>()
    .UseSqlite(connection)
    .UseAcmeTenancyForTesting(Services(tenants))
    .Options

ctx.Orders.Add(new Order { Total = 1 });      // TenantId bewusst leer
await ctx.SaveChangesAsync();
\end{sourcecode}

\code{Use\allowbreak{}Acme\allowbreak{}Tenancy\allowbreak{}For\allowbreak{}Testing(\allowbreak{}sp)}
bildet \code{Tenant\allowbreak{}Id} für jeden Provider als \code{Guid} ab (über einen
Value-Converter-Selector, wie EF es für Enums macht), ergänzt den SaveChanges-Interceptor des
Kerns und hinterlegt dafür die Mandanten-Property.
\code{Has\allowbreak{}Tenant\allowbreak{}Query\allowbreak{}Filters(\allowbreak{}this)} filtert
mit dem Mandanten als Abfrageparameter, derselbe Filter wie für Cosmos DB
(Kapitel~\ref{ch:without-sql-what-carries-over}). Row Level Security, Migrations-SQL, Session
und Translators fehlen. \code{Testing\allowbreak{}Mode\allowbreak{}Tests} betreibt den Modus
auf SQLite und InMemory.

\begin{important}[Ein Testmodus, kein Produktivmodus]
  \code{Use\allowbreak{}Acme\allowbreak{}Tenancy\allowbreak{}For\allowbreak{}Testing} lehnt
  PostgreSQL, SQL Server und Cosmos DB ab, damit er die echte Isolation nicht versehentlich
  ersetzt. Ob die Isolation hält, zeigen die Integrationstests des Plugins und die
  Integrationstests der Anwendung gegen die echte Datenbank (Testcontainers); der Testmodus
  hält nur die Tests der Geschäftslogik am Laufen.
\end{important}

\section{Was getestet wird}

\begin{datatable}{H{0.75}H{1.25}}
  \textbf{Test} & \textbf{Sichert ab} \\
  \tablehead
  Caching mit verschiedenen Optionen & Zwei Contexts mit verschiedenen Settings bekommen
  verschiedene Modelle (Hash-Fehler aus Kapitel~\ref{ch:the-options-extension}). \\
  \tablerule
  Ungültige Setting-Namen & Validierung inklusive Injection-Versuch. \\
  \tablerule
  Type Mapping ohne Konfiguration & \code{Tenant\allowbreak{}Id} wird \code{uuid};
  Literal-Format; anderer Spaltentyp wird nicht beansprucht. \\
  \tablerule
  \code{Current\allowbreak{}Tenant(\allowbreak{})} per
  \code{To\allowbreak{}Query\allowbreak{}String} & Übersetzung und dass nichts clientseitig
  ausgewertet wird (Kapitel~\ref{ch:preventing-client-side-evaluation}). \\
  \tablerule
  Explizit schlägt Convention & Eigene Policy gleichen Namens bleibt erhalten; Fluent API
  überstimmt das Attribut. \\
  \tablerule
  Warning-as-Error & Opt-out-Warnung wird mit \code{Configure\allowbreak{}Warnings} zum
  Fehler. \\
  \tablerule
  Validierung & RLS am abgeleiteten Typ und ungültige Policies werfen beim Modellbau. \\
  \tablerule
  Keine Phantom-Migration & Unverändertes Modell erzeugt keine einzige Operation. \\
  \tablerule
  Migrations-SQL & Initial, Änderung (DROP + CREATE), Opt-out, eigene Operationen. \\
  \tablerule
  SaveChanges-Logik & Setzen, fremder Mandant, Umhängen, kein Mandant. \\
  \tablerule
  Snapshot & Fluent API statt Annotations; Roundtrip ohne Unterschiede. \\
  \tablerule
  RLS in PostgreSQL & Isolation, keine Zeilen ohne Mandant, WITH CHECK greift. \\
  \tablerule
  Grenzfälle (Kapitel~\ref{ch:bulk-operations-and-the-limits}--\ref{ch:advanced-migrations})
  & ExecuteUpdate und rohes SQL gegen RLS, Value Generation, JSON-Spalten, Compiled Model,
  idempotentes Skript. \\
  \tablerule
  Kombinierbarkeit & Decorator um den Migrations-Generator und die Kompositionsmatrix zweier
  Plugins (Kapitel~\ref{ch:registering-services}), EFCore.NamingConventions in beiden
  Reihenfolgen (Kapitel~\ref{ch:conventions}). \\
  \tablerule
  Annotation-Format & Serialisierte Policies bleiben identisch; tolerantes Lesen
  (Kapitel~\ref{ch:versioning-and-compatibility}). \\
  \tablerule
  Context-Lebensdauern & DbContext-Factory, Hintergrundjobs, offene Verbindungen
  (Kapitel~\ref{ch:interceptors}, \ref{ch:using-the-plugin-in-asp}). \\
  \tablerule
  Mandantenübergreifender Zugriff & Reporting-Rolle liest alle Mandanten, kann aber nicht
  schreiben; die App-Rolle kann RLS nicht abschalten
  (Kapitel~\ref{ch:row-level-security-in-production}). \\
\end{datatable}

\section{Der Snapshot-Roundtrip}

Der wertvollste Test für Design-Time-Code: Snapshot generieren, mit Roslyn kompilieren,
Modell daraus bauen und gegen das Original diffen. Jede Abweichung zwischen Code-Generator
und Fluent API würde sonst bei jedem Nutzer eine Phantom-Migration erzeugen:

% source: tests/Acme.EntityFrameworkCore.Tenancy.Tests/SnapshotTests.cs
\begin{sourcecode}[tests/\allowbreak{}.\allowbreak{}.\allowbreak{}.\allowbreak{}/\allowbreak{}Snapshot\allowbreak{}Tests.\allowbreak{}cs (\allowbreak{}Auszug)]{csharp}
[Fact]
public void Snapshot_roundtrip_produces_no_differences()
{
    using var ctx = CreateContext();
    var snapshot = Compile(GenerateSnapshot(ctx));

    // Wie im EF-Migrator: finalisieren und initialisieren
    var model = snapshot.Model;
    if (model is IMutableModel mutable) model = mutable.FinalizeModel();
    model = ctx.GetService<IModelRuntimeInitializer>()
        .Initialize(model, designTime: true, validationLogger: null);

    var differences = ctx.GetService<IMigrationsModelDiffer>().GetDifferences(
        model.GetRelationalModel(), ctx.DesignModel().GetRelationalModel());

    Assert.Empty(differences);
}
\end{sourcecode}

\section{Integrationstests mit Testcontainers}

Das echte Sicherheitsverhalten lässt sich nur gegen PostgreSQL prüfen. Die Fixture startet
einen Container, legt das Schema mit \code{Ensure\allowbreak{}Created} an (inklusive RLS,
Kapitel~\ref{ch:how-migrations-work}), führt die eigenen Migrations-Operationen real aus und
verbindet die Tests mit einer eingeschränkten Rolle:

% source: tests/Acme.EntityFrameworkCore.Tenancy.IntegrationTests/PostgresFixture.cs
\begin{sourcecode}[tests/\allowbreak{}.\allowbreak{}.\allowbreak{}.\allowbreak{}Integration\allowbreak{}Tests/\allowbreak{}Postgres\allowbreak{}Fixture.\allowbreak{}cs (\allowbreak{}Auszug)]{csharp}
await owner.Database.EnsureCreatedAsync();          // RLS, Policies, Index

var commands = owner.GetService<IMigrationsSqlGenerator>().Generate(
[
    new CreateApplicationRoleOperation { Name = AppRole, CanLogin = true },
    new GrantTablePrivilegesOperation { Role = AppRole },
]);
await owner.GetService<IMigrationCommandExecutor>().ExecuteNonQueryAsync(
    commands, owner.GetService<IRelationalConnection>(),
    new MigrationExecutionState(), commitTransaction: true);
\end{sourcecode}

\begin{sourcecode}[tests/\allowbreak{}.\allowbreak{}.\allowbreak{}.\allowbreak{}Integration\allowbreak{}Tests/\allowbreak{}Row\allowbreak{}Level\allowbreak{}Security\allowbreak{}Tests.\allowbreak{}cs (\allowbreak{}Auszug, vereinfacht)]{csharp}
[Fact]
public async Task Database_rejects_rows_of_other_tenant()
{
    // Interceptor umgehen: direktes SQL. WITH CHECK der Policy muss greifen.
    var tenants = new MutableTenantProvider { CurrentTenantId = TenantId.New() };
    await using var db = _fixture.CreateAppContext(tenants);

    var ex = await Assert.ThrowsAsync<PostgresException>(() => db.Database.ExecuteSqlAsync(
        $"INSERT INTO \"Orders\" (\"TenantId\", \"Total\") VALUES ({Guid.NewGuid()}, 1)"));

    Assert.Equal(PostgresErrorCodes.InsufficientPrivilege, ex.SqlState);
}
\end{sourcecode}

\begin{tip}
  Für Provider und sehr tiefe Plugins stellt EF Core die \textbf{Specification Tests} bereit
  (\code{Microsoft.\allowbreak{}Entity\allowbreak{}Framework\allowbreak{}Core.\allowbreak{}Relational.\allowbreak{}Specification.\allowbreak{}Tests}).
  Das ist dieselbe Testsuite, gegen die EF und die Provider selbst getestet werden. Das
  NetTopologySuite-Plugin im EF-Repository zeigt, wie ein Plugin sie nutzt.
\end{tip}

\section{Kombinationen aus Migrationen und Isolation}

Einzelne Funktionen sind gut abgedeckt; die Fehler, die ein Review fand, lagen in Kombinationen:
ein geänderter Session-Schlüssel zusammen mit einer gelöschten Tabelle, eine Temporal-Tabelle,
eine ignorierte Mandanten-Property. Ein vollständiges kartesisches Produkt aller
Migrationsoperationen und Einstellungen wäre groß und meist redundant. Die Tests decken deshalb
eine kleine Matrix ab, ausgewählt nach einer Frage: Wo berühren zwei Funktionen dasselbe
Datenbankobjekt? Bei SQL Server sind das die gemeinsame Predicate-Funktion und die
schemagebundene Policy. Jeder Migrationsfall wird von EFs echtem Differ erzeugt, läuft Up und Down in
einer Transaktion gegen SQL Server und liest und schreibt danach als App-Login für zwei
Mandanten:

\begin{datatable}{H{1.05}H{0.95}H{1.00}}
  \textbf{Kombination} & \textbf{Test} & \textbf{Prüft} \\
  \tablehead
  Session-Schlüssel oder Policy-Schema geändert, eine geschützte Tabelle gelöscht &
  \code{Sql\allowbreak{}Server\allowbreak{}Session\allowbreak{}Settings\allowbreak{}Tests} &
  Reihenfolge des SQL, Up und Down, Isolation der verbleibenden Tabelle \\
  \tablerule
  Session-Schlüssel geändert, Tabelle umbenannt &
  \code{Sql\allowbreak{}Server\allowbreak{}Combined\allowbreak{}Migration\allowbreak{}Tests} &
  Reihenfolge des SQL, Up und Down, Policy auf dem neuen Namen \\
  \tablerule
  Policy-Schema geändert, Spalte umbenannt, Tabelle umbenannt und verschoben & wie oben & Up
  und Down, Funktion und Policy nur im neuen Schema \\
  \tablerule
  Migration scheitert, nachdem Policy und Funktion gelöscht wurden & wie oben & Rollback:
  Tabelle, Policy und Funktion mit dem alten Schlüssel sind unverändert \\
  \tablerule
  Temporal-Tabelle, ignorierte oder nicht gemappte Mandanten-Property & Modelltests beider
  relationaler Pakete & Beim Modellaufbau abgelehnt \\
\end{datatable}

Weiterhin ohne Test sind: Kombinationen bei PostgreSQL, wo Policies einer Umbenennung folgen und
es keine gemeinsame Funktion gibt, die Fehlerklasse also nicht entsteht; ein geänderter
Setting-Name des PostgreSQL-Pakets zusammen mit anderen Tabellenänderungen; mehrere Contexts,
die dieselbe Datenbank gleichzeitig migrieren; und das Cosmos-Paket, das keine Migrations hat.
Für diese Kombinationen wird nicht behauptet, dass sie funktionieren.

% ------------------------------------------------------------
\chapter{Versionierung und Kompatibilität}\label{ch:versioning-and-compatibility}
\versionrisk{medium}

\section{Paketversion und unterstützte EF-Version}

Das Beispiel versioniert Plugin und Handbuch gemeinsam: Die Paketversion entspricht der
Buchausgabe, Pakete 1.2.0 gehören zu Ausgabe~1.2. So gehört jeder Codeauszug zu genau einem
Paketstand. Welche EF-Version ein Stand unterstützt, steht deshalb nicht in der Versionsnummer,
sondern im Abhängigkeitsbereich der Pakete und in der Kompatibilitätstabelle der README: Plugin
1.x unterstützt EF Core~10. Der Wechsel auf eine neue EF-Major-Version ist für Anwender ein
Breaking Change und bekommt eine neue Major-Version, 2.0 für EF Core~11, zusammen mit einer
neuen Buchausgabe. Viele Plugins koppeln die Major-Version stattdessen an EF Core (10.x für EF
Core~10); das macht die passende Kombination an der Nummer sichtbar. Beide Wege funktionieren,
solange die Abhängigkeiten auf eine EF-Major-Version begrenzt sind. Eine Eigenschaft hält den
Bereich an einer Stelle, damit die CI ihn für den Patch-Kanal überschreiben kann (Abschnitt
„Patches und Previews in der CI“):

\begin{sourcecode}[Directory.\allowbreak{}Packages.\allowbreak{}props (\allowbreak{}Auszug)]{xml}
<PropertyGroup>
  <EFCoreVersion Condition="'$(EFCoreVersion)' == ''">[10.0.0,11.0.0)</EFCoreVersion>
  <CentralPackageFloatingVersionsEnabled>true</CentralPackageFloatingVersionsEnabled>
</PropertyGroup>
<PackageVersion Include="Microsoft.EntityFrameworkCore" Version="$(EFCoreVersion)" />
<PackageVersion Include="Npgsql.EntityFrameworkCore.PostgreSQL" Version="$(EFCoreVersion)" />
\end{sourcecode}

\section{Umgang mit interner API}

Plugins kommen um interne APIs selten herum: Das Beispiel erbt von
\code{Npgsql\allowbreak{}Annotation\allowbreak{}Provider} und
\code{Npgsql\allowbreak{}Annotation\allowbreak{}Code\allowbreak{}Generator} und nutzt
\code{Npgsql\allowbreak{}Options\allowbreak{}Extension}; insgesamt berührt das
PostgreSQL-Paket 10 interne Typen, das SQL-Server-Paket 2, das Cosmos-Paket 1, Kern und
relationale Schicht keinen (Inventar im Upgrade-Playbook). Solche Typen können sich auch in Minor-Versionen ändern. EF meldet ihre
Verwendung mit der Warnung \textbf{EF1001}.

\begin{datatable}{H{0.80}H{1.20}}
  \textbf{Maßnahme} & \textbf{Wirkung} \\
  \tablehead
  Öffentliche Alternative suchen & Jede vermiedene interne API ist ein Risiko weniger. \\
  \tablerule
  Interne Zugriffe bündeln & Wenige Klassen kapseln sie; ein Bruch betrifft eine Stelle. \\
  \tablerule
  \code{<\allowbreak{}No\allowbreak{}Warn>EF1001<\allowbreak{}/\allowbreak{}No\allowbreak{}Warn>}
  bewusst setzen & Projektweit und dokumentiert, nicht pro Zeile verstreut. \\
  \tablerule
  Versionsbereich begrenzen & Verhindert, dass Nutzer ungetestete Kombinationen bekommen. \\
  \tablerule
  CI gegen neueste Patches und Previews & Zeigt Brüche, bevor Anwender sie melden; die
  Kanäle beschreibt der Abschnitt „Patches und Previews in der CI“. \\
  \tablerule
  Roundtrip- und SQL-Tests & Bemerken Verhaltensänderungen, die der Compiler nicht sieht. \\
\end{datatable}

\begin{note}[Multi-Targeting]
  Mehrere EF-Major-Versionen gleichzeitig zu unterstützen ist per Multi-Targeting möglich,
  verdoppelt aber Testaufwand und bedingte Kompilierung. Für die meisten Plugins ist ein
  Release-Zweig pro EF-Major-Version einfacher.
\end{note}

\section{Annotation-Formate weiterentwickeln}

Alles, was ein Plugin in Annotations schreibt, landet in den Model-Snapshots seiner
Anwender, oft lange vor dem nächsten Upgrade. Serialisiert eine neue Plugin-Version dieselbe
Konfiguration anders (eine zusätzliche JSON-Property, eine andere Reihenfolge, eine andere
Enum-Schreibweise), vergleicht der Differ zwei verschiedene Strings. Jeder Anwender bekommt
dann eine Migration, die alle Policies ohne echte Änderung löscht und neu anlegt. Das
serialisierte Format gehört deshalb genauso zur öffentlichen API wie die Methoden.

\begin{datatable}{H{0.79}H{1.21}}
  \textbf{Regel} & \textbf{Wie} \\
  \tablehead
  Format per Test festschreiben & Ein Golden-Value-Test vergleicht die serialisierte Form
  Zeichen für Zeichen. \\
  \tablerule
  Tolerant lesen & Fehlende Properties werden zu Defaults, unbekannte werden ignoriert (das
  Standardverhalten von System.Text.Json). \\
  \tablerule
  Nur ergänzen, Defaults nicht schreiben & Neue optionale Properties mit
  \code{[Json\allowbreak{}Ignore(\allowbreak{}Condition =
  When\allowbreak{}Writing\allowbreak{}Default)]} lassen bestehende Annotations unverändert.
  \\
  \tablerule
  Echte Formatänderungen nur mit Major-Version & Dokumentieren, am besten zusammen mit dem
  EF-Major-Upgrade, wenn Anwender ohnehin mit Migrationen rechnen. \\
\end{datatable}

% source: tests/Acme.EntityFrameworkCore.Tenancy.Tests/AnnotationFormatTests.cs
\begin{sourcecode}[tests/\allowbreak{}.\allowbreak{}.\allowbreak{}.\allowbreak{}/\allowbreak{}Annotation\allowbreak{}Format\allowbreak{}Tests.\allowbreak{}cs (\allowbreak{}Auszug)]{csharp}
[Fact]
public void Serialized_policy_format_is_stable()
{
    var policy = new RlsPolicy("support_read", PolicyCommand.Select, Permissive: false,
        Roles: ["shop_support"], Using: "true", WithCheck: null);

    // Golden Value: eine Änderung bricht bestehende Snapshots
    Assert.Equal(
        """{"Name":"support_read","Command":"Select","Permissive":false"""
        + ""","Roles":["shop_support"],"Using":"true","WithCheck":null}""",
        policy.Serialize());
}
\end{sourcecode}

Zwei weitere Tests prüfen das tolerante Lesen in beide Richtungen. Der Snapshot-Roundtrip
aus Kapitel~\ref{ch:test-strategy} ergänzt sie: Der Roundtrip zeigt, dass der aktuelle
Code-Generator und die Fluent API übereinstimmen, der Golden Value zeigt, dass die nächste
Version noch zu den Snapshots passt, die Anwender bereits haben.

\section{Fehlerklassen beim Upgrade}

Ein Upgrade auf eine neue EF-Major-Version bricht ein Plugin selten nur an einer Stelle,
und nicht jeder Bruch meldet sich. Wer die Fehlerklassen kennt, weiß, welcher Test sie
findet, und vor allem, welche er nicht findet:

\begin{datatable}{H{0.70}H{1.10}H{1.20}}
  \textbf{Fehlerklasse} & \textbf{Typisches Beispiel} & \textbf{Im Beispiel erkannt durch}
  \\
  \tablehead
  Compile-Time Break & Konstruktor einer Provider-Klasse hat einen neuen Parameter, ein
  interner Typ ist umgezogen & Build; \code{Public\allowbreak{}API}-Dateien für die eigene
  API (Kapitel~\ref{ch:ci-and-api-stability}) \\
  \tablerule
  Runtime Break & Ein Dienst wechselt die Lebensdauer, eine Registrierung schlägt fehl &
  \code{Service\allowbreak{}Registration\allowbreak{}Tests},
  \code{Design\allowbreak{}Time\allowbreak{}Services\allowbreak{}Tests},
  \code{Options\allowbreak{}Tests}, \code{Composition\allowbreak{}Tests} \\
  \tablerule
  Silent Behavioral Change & Eine Reihenfolge, ein Vorrang oder eine Nullability-Regel von EF
  ändert sich & \code{Failure\allowbreak{}Case\allowbreak{}Tests},
  \code{Composition\allowbreak{}Matrix\allowbreak{}Tests},
  \code{State\allowbreak{}And\allowbreak{}Lifetime\allowbreak{}Tests},
  \code{Save\allowbreak{}Changes\allowbreak{}Interceptor\allowbreak{}Tests}; mit Datenbank
  \code{Materialization\allowbreak{}Tests}, \code{Connection\allowbreak{}Lifetime\allowbreak{}Tests} \\
  \tablerule
  SQL Change & Der Provider schreibt einen Ausdruck anders oder optimiert eine Bedingung weg &
  \code{Query\allowbreak{}Translation\allowbreak{}Tests},
  \code{Deep\allowbreak{}Query\allowbreak{}Tests}, \code{Debugging\allowbreak{}Tests} \\
  \tablerule
  Model Change & Eine Convention läuft früher, eine Standardkonfiguration ändert sich &
  \code{Convention\allowbreak{}Tests}, \code{Tenant\allowbreak{}Property\allowbreak{}Tests},
  \code{Validation\allowbreak{}Tests}, \code{Type\allowbreak{}Mapping\allowbreak{}Tests},
  \code{Json\allowbreak{}Tests}, \code{Value\allowbreak{}Generation\allowbreak{}Tests},
  \code{Naming\allowbreak{}Convention\allowbreak{}Tests} \\
  \tablerule
  Snapshot/Migration Change & Der Code-Generator oder der Differ verhält sich anders; jeder
  Anwender bekäme eine Phantom-Migration & \code{Snapshot\allowbreak{}Tests} (Roundtrip),
  \code{Migration\allowbreak{}Sql\allowbreak{}Tests},
  \code{Migration\allowbreak{}Script\allowbreak{}Tests},
  \code{Annotation\allowbreak{}Format\allowbreak{}Tests},
  \code{Compiled\allowbreak{}Model\allowbreak{}Tests}, Verbrauchertest \\
  \tablerule
  Performance Regression & Ein heißer Pfad allokiert mehr oder wird langsamer & Das
  Benchmark-Projekt (Kapitel~\ref{ch:plugin-performance}) \\
\end{datatable}

Die dritte Zeile ist die gefährlichste. Viele Tests des Beispiels prüfen nicht das Plugin,
sondern das Verhalten von EF, auf das sich das Plugin und dieses Buch verlassen. Sie sind
beim Upgrade wertvoller als jeder andere Test, weil sie genau das melden, was kein Compiler
sieht.

\section{Das Upgrade-Playbook}

Der Ablauf ist am Wechsel von EF Core~10 auf EF Core~11 beschrieben, gilt aber für jede
Major-Version: Die Nummern sind austauschbar, die Reihenfolge nicht. Vorher wird die
bisherige Linie in einem Release-Zweig eingefroren; sie bekommt weiter Patches für ihre
EF-Version, während \code{main} auf die neue wechselt.

\begin{sourcecode}[Schritt~1: Zweig und Versionen]{text}
git switch -c release/1.x main && git push -u origin release/1.x
git switch main

global.json               "version": "11.0.100"
Directory.Build.props     <TargetFramework>net11.0</TargetFramework>
                          <VersionPrefix>2.0.0</VersionPrefix>
Directory.Packages.props  <EFCoreVersion ...>[11.0.0,12.0.0)</EFCoreVersion>
\end{sourcecode}

Welches Target-Framework eine EF-Version voraussetzt, folgt keiner festen Regel: EF Core~9
zielt auf \code{net8.\allowbreak{}0}, EF Core~10 auf \code{net10.\allowbreak{}0}, die
Release Candidates von EF Core~11 auf \code{net11.\allowbreak{}0}
(\href{https://learn.microsoft.com/en-us/ef/core/what-is-new/}{EF Core releases and planning}). Vor dem Upgrade gehört deshalb ein Blick in
die Pakete von EF \emph{und} Provider: Welche Frameworks bieten sie an, welches SDK
brauchen sie? Im Beispiel verlangt EF Core~11 \code{net11.\allowbreak{}0}, also wechseln
Target-Framework, SDK und Pakete gemeinsam; ein neues EF-Paket auf dem alten
Target-Framework ließe sich nicht einmal wiederherstellen.

\begin{datatable}{H{0.24}H{1.27}H{1.76}H{0.73}}
  \textbf{Nr.} & \textbf{Schritt} & \textbf{Vorgehen} & \textbf{Im Beispiel} \\
  \tablehead
  1 & Pakete aktualisieren & Versionen aus Schritt~1, dazu Begleitpakete, die mit EF erscheinen
  (\code{Microsoft.\allowbreak{}Extensions.\allowbreak{}*}); fest gepinnte Begleitpakete
  blockieren schon den Restore & \code{EFCore\allowbreak{}Version} \\
  \tablerule
  2 & Compilerfehler klassifizieren & Nach Fehlercode zählen: verschobene Typen
  (CS0234, CS0246), geänderte Konstruktoren und Dependencies (CS1729, CS7036), entfallene
  Überschreibungen (CS0115) & Befehl unten \\
  \tablerule
  3 & Interne API untersuchen & Das EF1001-Inventar (unten) gegen die Quellen beider
  Versionen vergleichen; jede Stelle einzeln, auch die, die noch kompilieren;
  mit den volatilen Kapiteln beginnen (Seite~\pageref{sec:version-dependency}) &
  Kern: 0, PostgreSQL: 10, SQL Server: 2, Cosmos DB: 1 interne Typen (Funktionen-Beispiel: 5) \\
  \tablerule
  4 & Registrierungen und Lebensdauern vergleichen & EFs Service-Tabellen beider Versionen
  vergleichen; die eigenen \code{Replace}-Aufrufe legen Lebensdauern fest &
  \code{Service\allowbreak{}Registration\allowbreak{}Tests} \\
  \tablerule
  5 & Abgeleitete Provider-Klassen prüfen & Der Compiler meldet entfallene Methoden, aber
  nicht neue virtuelle Überladungen, die die Basis jetzt stattdessen aufruft: Basisklasse
  im Diff lesen & \code{Migration\allowbreak{}Sql\allowbreak{}Tests},
  \code{Deep\allowbreak{}Query\allowbreak{}Tests} \\
  \tablerule
  6 & Modell und Conventions prüfen & Modell-Tests; zusätzlich
  \code{To\allowbreak{}Debug\allowbreak{}String} des Design-Time-Modells beider Versionen
  vergleichen (Anhang~\ref{app:D}) & \code{Convention\allowbreak{}Tests} u.\,a. \\
  \tablerule
  7 & Annotation-Formate prüfen & Das eigene Format darf sich nicht ändern; EFs eigene
  Annotations, die das Plugin liest, können es & \code{Annotation\allowbreak{}Format\allowbreak{}Tests} \\
  \tablerule
  8 & Query-SQL testen & Alle \code{To\allowbreak{}Query\allowbreak{}String}-Tests; jede
  Abweichung erklären, nicht nur die Erwartung anpassen & \code{Query\allowbreak{}Translation\allowbreak{}Tests} u.\,a. \\
  \tablerule
  9 & Migrations-SQL vergleichen & Generiertes SQL der Tests; bei Beispielanwendungen das
  Skript von \code{dotnet ef migrations script} vor und nach dem Upgrade vergleichen &
  \code{Migration\allowbreak{}Sql\allowbreak{}Tests},
  \code{Migration\allowbreak{}Script\allowbreak{}Tests} \\
  \tablerule
  10 & Snapshot-Roundtrip & Snapshot erzeugen, kompilieren, vergleichen; eine neue Migration
  auf einem bestehenden Projekt muss leer sein & \code{Snapshot\allowbreak{}Tests} \\
  \tablerule
  11 & Compiled Models & Generieren, kompilieren, mit \code{Use\allowbreak{}Model} laden, auch die
  Laufzeit-Annotations & \code{Compiled\allowbreak{}Model\allowbreak{}Tests},
  \code{run.\allowbreak{}sh} (Seite~\pageref{sec:consumer-test}) \\
  \tablerule
  12 & Trimming und NativeAOT & Neue IL-Warnungen (IL2xxx, IL3xxx) im Build lesen und eine
  NativeAOT-Anwendung starten; die experimentelle API für vorkompilierte Abfragen (EF9100)
  ändert sich besonders oft & \code{Is\allowbreak{}Aot\allowbreak{}Compatible},
  \code{run.\allowbreak{}sh -\allowbreak{}-nativeaot} \\
  \tablerule
  13 & Integrationstests & Gegen echte Datenbanken, mit den Provider-Versionen der neuen
  Linie & Integrations- und SQL-Server-Tests \\
  \tablerule
  14 & Eigene öffentliche API vergleichen & Der Public-API-Analyzer bricht den Build bei jeder
  Abweichung; gewollte Brüche bewusst eintragen & \code{Public\allowbreak{}API.\allowbreak{}*.\allowbreak{}txt} \\
  \tablerule
  15 & Benchmarks vergleichen & Dasselbe Benchmark auf beiden Linien; Verhältnis und
  Allokationen vergleichen, nicht absolute Zeiten & Benchmark-Projekt \\
  \tablerule
  16 & Breaking Changes dokumentieren & CHANGELOG, Kompatibilitätstabelle (Plugin~2.x zu
  EF~Core~11), Hinweise für Anwender, Supportdauer der alten Linie & \code{CHANGELOG.\allowbreak{}md} \\
\end{datatable}

Für die Schritte 2 und 3 genügen zwei Befehle. Der zweite hebt die projektweite
Unterdrückung von EF1001 nur für diesen Build auf und zeigt so jede Stelle, an der das
Plugin interne API berührt:

\begin{sourcecode}[Schritte~2 und 3]{text}
# Compilerfehler nach Code
dotnet build -c Release 2>&1 | grep -oE "error [A-Z]+[0-9]+" | sort | uniq -c | sort -rn

# Inventar interner API: NoWarn ohne EF1001 (EF9100 bleibt unterdrückt)
dotnet build -c Release --no-incremental -p:NoWarn=EF9100 2>&1 \
  | grep -oE "EF1001: [A-Za-z.]+" | sort | uniq -c | sort -rn
\end{sourcecode}

\phantomsection\label{sec:version-dependency}
Wie stark ein Kapitel von der EF-Version abhängt, steht an seinem Anfang. Bei einem
Major-Upgrade prüft man die \textbf{volatilen} Kapitel zuerst: Ihr Code leitet von Klassen
aus \code{.Internal}-Namespaces ab oder nutzt sie direkt (EF1001), und Brüche dort zeigen
sich oft erst zur Laufzeit. Danach folgen die mittleren, die stabilen zuletzt.

\begin{datatable}{H{0.55}H{1.15}H{1.30}}
  \textbf{Einstufung} & \textbf{Bedeutung} & \textbf{Kapitel} \\
  \tablehead
  \textcolor{risk@volatile}{\textbf{volatil}} & Stark von \code{.Internal}-Namespaces
  abhängig (EF1001); bei jedem Major-Upgrade zuerst prüfen & \ref{ch:declarative-sql-from-annotations}, \ref{ch:design-time-services}, \ref{ch:custom-sql-expression-nodes}, \ref{ch:property-level-fluent-api-and}, \ref{ch:scaffolding-in-practice} \\
  \tablerule
  \textcolor{risk@medium}{\textbf{mittel}} & Provider-Basisklassen oder öffentliche, aber oft
  geänderte API & \ref{ch:ef-core-s-architecture-from}, \ref{ch:the-options-extension}, \ref{ch:registering-services}, \ref{ch:type-mappings-and-value-converters}, \ref{ch:the-query-pipeline}, \ref{ch:translating-methods-and-members}, \ref{ch:further-query-extension-points}, \ref{ch:custom-migration-operations}, \ref{ch:compiled-models-trimming-and-nativeaot}, \ref{ch:scaffolding-reverse-engineering}, \ref{ch:test-strategy}, \ref{ch:versioning-and-compatibility}, \ref{ch:bulk-operations-and-the-limits}, \ref{ch:the-update-pipeline}, \ref{ch:complex-types-and-json}, \ref{ch:aggregate-translators}, \ref{ch:a-second-provider-sql-server},
  \ref{ch:without-sql-what-carries-over} \\
  \tablerule
  \textcolor{risk@stable}{\textbf{stabil}} & Konzepte und öffentliche API & alle übrigen \\
\end{datatable}

\begin{tip}[Abweichungen erklären, nicht wegtesten]
  Schlägt nach dem Upgrade ein SQL- oder Verhaltenstest fehl, ist die Versuchung groß, nur
  die Erwartung anzupassen. Jede Abweichung gehört aber erst erklärt: Ist sie eine
  Verbesserung von EF, ein Bruch des Plugins oder eine geänderte Semantik, die Anwender
  betrifft? Erst die Antwort entscheidet, ob der Test, der Code oder die Dokumentation
  geändert wird.
\end{tip}

\section{Patches und Previews in der CI}

Die Versionsbereiche lösen nur einen Teil des Problems. NuGet nimmt aus einem Bereich die
\emph{niedrigste} Version, die alle Abhängigkeiten erfüllt; die reguläre CI testet also
gegen den Anfang der Linie, während Anwender neuere Patches bekommen. Die CI des Beispiels
arbeitet deshalb mit drei Kanälen:

\begin{datatable}{H{0.87}H{1.27}H{1.07}H{0.79}}
  \textbf{Kanal} & \textbf{Auslöser und Pakete} & \textbf{Zweck} & \textbf{Bei Fehler} \\
  \tablehead
  Regulär (\code{ci.\allowbreak{}yml}) & Jeder Push und Pull Request; niedrigste Version des
  Bereichs & Die Linie, gegen die das Plugin gebaut wird & Pflicht-Check \\
  \tablerule
  Patches (\code{compatibility.\allowbreak{}yml}) & Wöchentlich und manuell;
  \code{-p:EFCore\allowbreak{}Version=10.0.*}, nur stabile Pakete & Neueste Patches der
  eigenen Linie & Untersuchen, Patch-Release \\
  \tablerule
  Nächste Major-Version & Eigener Job, zeitgesteuert; Previews und RCs & Brüche Monate vor
  dem Release sehen & Kein Pflicht-Check \\
\end{datatable}

Der Patch-Kanal ist im Beispiel aktiv. Beim ersten Lauf zeigte er bereits einen echten
Befund: Das Plugin selbst lief mit EF Core~10.0.14, aber Test- und Beispielprojekte ließen
sich nicht wiederherstellen, weil sie Begleitpakete fest auf 10.0.0 und eine eigene
Roslyn-Version gepinnt hatten. Seitdem folgen diese Pakete \code{EFCore\allowbreak{}Version}.

Ein Kanal für die nächste Major-Version ist im Beispiel bewusst nicht aktiv: Er arbeitet mit
Vorabversionen. Wer ihn einrichtet, hält vier Regeln ein, damit normale Releases nie an
Preview-Pakete gekoppelt werden:

\begin{datatable}{H{0.85}H{1.15}}
  \textbf{Regel} & \textbf{Warum} \\
  \tablehead
  Alles über globale Eigenschaften dieses Jobs umstellen: SDK, Target-Framework und
  \code{EFCore\allowbreak{}Version} & Die Standardwerte in den Projektdateien bleiben
  stabil; ein Preview-Paket kann so nie in einen regulären Build gelangen \\
  \tablerule
  SDK, Target-Framework und Pakete gemeinsam wechseln & Mit dem alten Target-Framework
  scheitert schon der Restore; ein solcher Job kann nie grün werden \\
  \tablerule
  Nie packen oder veröffentlichen & Releases entstehen nur aus regulären Builds der
  Release-Zweige \\
  \tablerule
  Kein Pflicht-Check & Previews brechen regelmäßig; der Job ist eine Frühwarnung, die ein
  Issue erzeugt, kein Tor für Pull Requests \\
\end{datatable}

\begin{sourcecode}[Job für die nächste EF-Version (\allowbreak{}Skizze,\allowbreak{} nicht im Repository)]{text}
  ef-next:
    runs-on: ubuntu-latest
    continue-on-error: true                 # Frühwarnung, kein Pflicht-Check
    steps:
      - uses: actions/checkout@v4
      - uses: actions/setup-dotnet@v4
        with: { dotnet-version: 11.0.x, dotnet-quality: preview }
      - run: dotnet new globaljson --sdk-version "$(dotnet --version)" --force
      - name: Build against the next EF Core version (never packed or published)
        run: dotnet build -c Release -p:TargetFramework=net11.0 -p:EFCoreVersion=11.0.0-*
      - run: >
          dotnet test tests/Acme.EntityFrameworkCore.Tenancy.Tests -c Release --no-build
          -p:TargetFramework=net11.0
\end{sourcecode}

Test-Abhängigkeiten von Drittanbietern wie EFCore.NamingConventions erscheinen oft erst
später für die neue Version und können den Restore dieses Jobs blockieren. Solche Tests
nimmt man im Preview-Kanal gezielt aus, statt den Kanal abzuschalten.

\section{NuGet-Advisories und transitive Abhängigkeiten}\label{sec:nuget-advisories}

Ein Plugin liefert nicht nur eigenen Code aus, sondern auch seinen Abhängigkeitsgraphen. Das
NuGet-Audit meldet bekannte Schwachstellen als Warnungen NU1901 bis NU1904
(\href{https://learn.microsoft.com/en-us/nuget/concepts/auditing-packages}{NuGet: Auditing package dependencies}).
Entscheidend ist, \emph{wen} ein Befund trifft. Ermittelt wird das mit \code{dotnet nuget why},
das den Pfad von der direkten Abhängigkeit bis zum betroffenen Paket zeigt. Im Beispiel gab es
drei Fälle:

\begin{datatable}{H{0.75}H{1.05}H{0.70}H{1.50}}
  \textbf{Paket} & \textbf{Herkunft} & \textbf{Betrifft} & \textbf{Entscheidung} \\
  \tablehead
  \code{System.\allowbreak{}Security.\allowbreak{}Cryptography.\allowbreak{}Pkcs} 9.0.4 &
  EF Core SqlServer → \code{Microsoft.\allowbreak{}Data.\allowbreak{}SqlClient} 6.1 & Verbraucher
  des SQL-Server-Pakets & Kein Patch und keine SqlClient-Version (auch nicht 7.x) bringt bisher
  9.0.21. Deshalb eine direkte Abhängigkeit auf 9.0.22 im SQL-Server-Paket; sie landet im
  nuspec und korrigiert den Graphen der Anwendung. \\
  \tablerule
  \code{Cryptography.\allowbreak{}Xml}/\code{Pkcs} 9.0.0 & EF Core Design 10.0.0 → MSBuild
  17.14, Workspaces.MSBuild 4.14 & Nur Werkzeuge: Tests, Beispiel & Neuere Design-Patches
  brauchen andere MSBuild-Pakete, würden aber EF selbst mitziehen und den Test der
  Mindestversion aufheben. Deshalb \code{Xml} 9.0.22 nur in Test- und Beispielprojekten. \\
  \tablerule
  \code{SSH.\allowbreak{}NET} 2024.2.0 & Testcontainers 4.6 & Nur Tests & Testcontainers 4.16
  innerhalb derselben Major-Version; bringt 2026.0.0. \\
\end{datatable}

Jede direkte Abhängigkeit auf ein eigentlich transitives Paket steht in
\code{Directory.\allowbreak{}Packages.\allowbreak{}props} mit Herkunft, Advisory und dem
Kriterium, wann sie wieder entfallen kann. Unterdrückt wird nichts: Weder \code{NoWarn} noch ein
abgeschaltetes Audit machen eine Abhängigkeit sicher, sie verstecken nur den Befund. Beide
CI-Kanäle führen zusätzlich \code{dotnet list package -\allowbreak{}-vulnerable
-\allowbreak{}-include-transitive} aus und scheitern an jedem Fund. Der reguläre Lauf baut
weiterhin gegen EF Core 10.0.0, der Kompatibilitätslauf gegen den neuesten 10.0.x-Patch; die
Korrekturen sind Mindestversionen und gelten in beiden. Für Werkzeuge gibt es eine Wechselwirkung:
\code{Microsoft.\allowbreak{}Entity\allowbreak{}Framework\allowbreak{}Core.\allowbreak{}Tasks}
wertet die Audit-Warnungen beim Veröffentlichen als Fehler, sodass NativeAOT mit EF Design
10.0.0 scheitert (Seite~\pageref{sec:consumer-test}).

\section{Checkliste für das Upgrade}

\begin{checklist}
  \item Release-Zweig für die alte Linie anlegen; \code{main} auf die neue Version heben.
  \item SDK, Target-Framework, \code{EFCore\allowbreak{}Version}, Begleitpakete und
    \code{Version\allowbreak{}Prefix} gemeinsam ändern.
  \item Compilerfehler nach Code zählen und klassenweise abarbeiten.
  \item EF1001-Inventar erstellen und jede Stelle gegen die neuen Quellen prüfen, auch die
    kompilierenden.
  \item Lebensdauern gegen EFs Service-Tabelle prüfen; abgeleitete Provider-Klassen auf neue
    virtuelle Überladungen lesen.
  \item Modell-, SQL-, Migrations-, Snapshot- und Compiled-Model-Tests grün; jede Abweichung
    erklärt.
  \item Annotation-Formate unverändert (Golden-Value-Tests).
  \item Neue IL-Warnungen für Trimming und AOT gelesen.
  \item Integrationstests gegen echte Datenbanken mit den neuen Providern.
  \item Public-API-Diff geprüft; gewollte Brüche eingetragen.
  \item Benchmarks beider Linien verglichen.
  \item CHANGELOG, Kompatibilitätstabelle und Supportdauer der alten Linie dokumentiert.
\end{checklist}

% ------------------------------------------------------------
\chapter{Packaging und Dokumentation}\label{ch:packaging-and-documentation}
\versionrisk{stable}

Gemeinsame Einstellungen liegen in \code{Directory.\allowbreak{}Build.\allowbreak{}props},
Paketversionen zentral in \code{Directory.\allowbreak{}Packages.\allowbreak{}props} (Central
Package Management):

\begin{sourcecode}[Directory.\allowbreak{}Build.\allowbreak{}props (\allowbreak{}gekürzt)]{xml}
<PropertyGroup>
  <TargetFramework>net10.0</TargetFramework>
  <Nullable>enable</Nullable>
  <ManagePackageVersionsCentrally>true</ManagePackageVersionsCentrally>
  <NoWarn>$(NoWarn);EF1001</NoWarn>
  <Deterministic>true</Deterministic>
  <WarningsAsErrors>$(WarningsAsErrors);RS0016;RS0017</WarningsAsErrors>
</PropertyGroup>
<PropertyGroup>
  <PackageLicenseExpression>MIT</PackageLicenseExpression>
  <PackageReadmeFile>README.md</PackageReadmeFile>
  <VersionPrefix>1.2.0</VersionPrefix>       <!-- = Buchausgabe -->
  <IncludeSymbols>true</IncludeSymbols>
  <SymbolPackageFormat>snupkg</SymbolPackageFormat>
  <PublishRepositoryUrl>true</PublishRepositoryUrl>      <!-- SourceLink -->
  <EmbedUntrackedSources>true</EmbedUntrackedSources>
</PropertyGroup>
\end{sourcecode}

\begin{datatable}{H{0.70}H{1.30}}
  \textbf{Element} & \textbf{Warum} \\
  \tablehead
  \code{Generate\allowbreak{}Documentation\allowbreak{}File} & XML-Doku im Paket: Nutzer
  sehen Erklärungen in der IDE. \\
  \tablerule
  \code{Warnings\allowbreak{}As\allowbreak{}Errors} RS0016, RS0017 & Eine geänderte
  öffentliche API bricht den Build (Kapitel~\ref{ch:ci-and-api-stability}). \\
  \tablerule
  \code{Is\allowbreak{}Aot\allowbreak{}Compatible} & Analyzer melden AOT- und
  Trimming-Probleme beim Build. \\
  \tablerule
  SourceLink, \code{.\allowbreak{}snupkg} & Nutzer können in den Plugin-Code debuggen. Seit
  dem .NET-8-SDK eingebaut. \\
  \tablerule
  README im Paket & Erscheint auf nuget.org; Minimalbeispiel und Kompatibilität. \\
  \tablerule
  CHANGELOG & Breaking Changes klar markieren, besonders bei EF-Major-Wechseln. \\
  \tablerule
  \code{Internals\allowbreak{}Visible\allowbreak{}To} & Tests erreichen interne Logik ohne
  öffentliche API. \\
  \tablerule
  Paketnamen & Kern ohne Provider, Provider-Paket mit Suffix
  (\code{.\allowbreak{}Postgre\allowbreak{}SQL},
  \code{.\allowbreak{}Sql\allowbreak{}Server}). \\
\end{datatable}

Zur Dokumentation gehören neben der API-Doku vor allem die Dinge, die Nutzer nicht im Code
sehen: welche Einzeldienste ersetzt werden (Kombinierbarkeit), welche Event-IDs es gibt,
welche Einschränkungen bestehen (Pooling) und welche Sicherheitsannahmen gelten
(Datenbankrollen).

% ------------------------------------------------------------
\chapter{Die Checkliste von A bis Z}\label{ch:the-checklist-from-a-to}
\versionrisk{stable}

Für jedes neue Plugin, in der Reihenfolge der Entwicklung. Jeder Punkt fasst eine Regel zusammen; die Begründung steht in den Kapiteln, die unter jeder Überschrift genannt sind.

\section{Architektur}

\noindent Details: Kapitel~\ref{ch:the-sample-scenario}, \ref{ch:choosing-the-right-extension-point}.

\begin{checklist}
  \item Stufe festlegen: Hilfsmethoden, Modell-Helfer, echtes Plugin oder Provider?
  \item Kern- und Provider-Paket trennen; Typen für Entitäts-Assemblies in den Kern.
  \item Für jeden Erweiterungspunkt zuerst ein Plugin-Interface suchen, erst dann ersetzen.
  \item Jeden Mechanismus mit der Entscheidungshilfe gegen die einfachere Alternative
    prüfen.
\end{checklist}

\section{Options}

\noindent Details: Kapitel~\ref{ch:the-options-extension}.

\begin{checklist}
  \item Extension immutable mit Kopierkonstruktor und virtuellem Clone.
  \item Alle modellrelevanten Optionen im Hash, nichts, was pro Context variiert.
  \item Validate für Voraussetzungen (Provider, Kombinationen).
  \item LogFragment und PopulateDebugInfo pflegen.
  \item Use-Methode im Namespace Microsoft.EntityFrameworkCore, generisch und
    nicht-generisch.
\end{checklist}

\section{Dienste}

\noindent Details: Kapitel~\ref{ch:registering-services}, \ref{ch:state-lifetimes-and-thread-safety}.

\begin{checklist}
  \item Registrierung über den ServicesBuilder (Lebensdauern aus EFs Tabelle).
  \item Einzeldienste dekorieren, wo das öffentliche Interface genügt.
  \item Sonst Replace, und die Logik öffentlich und virtuell anbieten.
  \item ISingletonOptions für Singleton-Dienste, doppelt registriert.
  \item Jeden Zustand einordnen: Anfragedaten nie in
    Options, Hash, Modell, Singleton oder übersetzte Abfrage.
  \item Scoped-Dienste mit Zustand für Pooling zurücksetzen (IResettableService).
  \item App-Abhängigkeiten über AddInterceptors, Pooling-Verhalten dokumentieren.
\end{checklist}

\section{Modell}

\noindent Details: Kapitel~\ref{ch:annotations-and-fluent-api}, \ref{ch:conventions}, \ref{ch:model-validation}.

\begin{checklist}
  \item Annotations mit Präfix, nur einfache Werte, keine Defaults speichern.
  \item Fluent API nicht-generisch implementieren, generisch delegieren.
  \item Convention-Builder-Überladungen für ConfigurationSource.
  \item Conventions bei ModelFinalizing, wenn sie vom fertigen Modell abhängen.
  \item Attribute im Kernpaket, Attribut-Conventions im Provider-Paket.
  \item Validierung als IModelFinalizedConvention: Fehler werfen, Fragwürdiges warnen.
  \item Tabellenabbildungen gegen die Unterstützungsmatrix prüfen und nicht unterstützte beim
    Modellbau ablehnen (Seite~\pageref{sec:supported-table-mappings}).
\end{checklist}

\section{Type Mappings}

\noindent Details: Kapitel~\ref{ch:type-mappings-and-value-converters}.

\begin{checklist}
  \item Geschützter Parameter-Konstruktor und Clone-Überschreibung.
  \item Literal-Erzeugung für Modell- und Provider-Typ.
  \item Öffentliche statische Default-Property, gemeinsame Konverter-Instanz.
  \item Gemeinsame Value-Comparer-Instanzen, sonst erzeugt EF sie per Reflection (NativeAOT).
  \item JsonValueReaderWriter mit Instance und ConstructorExpression.
\end{checklist}

\section{Abfragen}

\noindent Details: Kapitel~\ref{ch:the-query-pipeline}, \ref{ch:translating-methods-and-members}, \ref{ch:preventing-client-side-evaluation}.

\begin{checklist}
  \item EF.Functions-Methoden im Kern, Übersetzung im Provider-Paket.
  \item Translators: null für fremde Methoden, Nullability korrekt, Type Mappings setzen.
  \item IEvaluatableExpressionFilterPlugin für Methoden ohne Datenbezug.
  \item Keine Laufzeitwerte zur Übersetzungszeit einsetzen (Query-Cache, AOT).
\end{checklist}

\section{Laufzeit}

\noindent Details: Kapitel~\ref{ch:interceptors}, \ref{ch:diagnostics-and-logging}.

\begin{checklist}
  \item Interceptors synchron und asynchron implementieren.
  \item Tracking-Markierung von tatsächlicher Wertänderung unterscheiden
    (Seite~\pageref{sec:tracking-flag-versus-value}).
  \item Singleton-Interceptors zustandslos und threadsicher.
  \item Singleton-Interceptors in ApplyServices registrieren oder als eine geteilte Instanz
    übergeben, nie pro Context neu; Anfragedaten pro Aufruf über den Context holen.
  \item Eigene Event-IDs in eigenem Bereich, EventDefinitions gecacht.
  \item Sicherheitsrelevantes als Warnung, eskalierbar per ConfigureWarnings.
  \item Wo Interceptors nicht greifen, eine explizite API anbieten (z. B. ApplyTenantSession).
\end{checklist}

\section{Migrations \& Design-Time}

\noindent Details: Kapitel~\ref{ch:how-migrations-work}, \ref{ch:declarative-sql-from-annotations}, \ref{ch:custom-migration-operations}, \ref{ch:design-time-services}, \ref{ch:compiled-models-trimming-and-nativeaot}.

\begin{checklist}
  \item Deklarativ über Annotation-Provider (nur Design-Time) und SQL-Generator.
  \item Bezeichner maskieren, Literale über Type Mappings, selbst erzeugtes SQL validieren.
  \item Eigene Operationen idempotent, Extensions im Migrations-Namespace.
  \item Schemagebundene Objekte bei Umbenennungen vor der Änderung entfernen und danach mit
    den endgültigen Namen neu anlegen (Seite~\pageref{sec:sqlserver-schema-bound-changes}).
  \item DesignTimeServicesReference mit forProvider im Startprojekt der Anwendung (nicht im Plugin); Code-Generator entfernt Annotations.
  \item Compiled Models prüfen: Default, Instance, Laufzeit-Annotations.
\end{checklist}

\section{Qualität}

\noindent Details: Kapitel~\ref{ch:test-strategy}, \ref{ch:versioning-and-compatibility}, \ref{ch:packaging-and-documentation}.

\begin{checklist}
  \item Modell-, SQL-, Snapshot-Roundtrip- und Integrationstests.
  \item Caching-Test mit verschiedenen Optionen; Test auf keine Phantom-Migration.
  \item Versionsbereich je EF-Major, EF1001 bewusst; CI mit Patch-Kanal, Preview-Kanal nur als
    Frühwarnung; Upgrades nach dem Playbook.
  \item XML-Doku, SourceLink, Symbole, README, CHANGELOG, Einschränkungen dokumentiert.
  \item Die Pakete als Verbraucher prüfen: Tooling, Compiled Model, NativeAOT
    (Seite~\pageref{sec:consumer-test}).
  \item Serialisierte Annotation-Formate mit Golden-Value-Tests festschreiben.
  \item Zusammen mit verbreiteten Plugins testen (Naming Conventions, Decorators) und die
    Koexistenz mit einem zweiten Plugin in beiden Reihenfolgen prüfen
    (Kapitel~\ref{ch:registering-services}).
\end{checklist}

\section{Fortgeschrittene Erweiterungspunkte (Teil~\ref{part:advanced})}

\noindent Details: Kapitel~\ref{ch:bulk-operations-and-the-limits}, \ref{ch:value-generation}, \ref{ch:complex-types-and-json}, \ref{ch:advanced-migrations}.

\begin{checklist}
  \item Wege dokumentieren, die Interceptors umgehen (ExecuteUpdate, rohes SQL).
  \item Value Generators mit HasValueGeneratorFactory(Type) konfigurieren (Compiled Models).
  \item Conventions für Complex Types und JSON-Spalten prüfen.
  \item Reihenfolge eigener Migrations-Operationen dokumentieren und testen.
  \item Idempotente Skripte und suppressTransaction berücksichtigen.
\end{checklist}

\section{Betrieb und Integration (Teil~\ref{part:operations})}

\noindent Details: Kapitel~\ref{ch:row-level-security-in-production}, \ref{ch:tenant-safe-relationships}, \ref{ch:retrofitting-and-operating}, \ref{ch:plugin-performance}, \ref{ch:using-the-plugin-in-asp}, \ref{ch:ci-and-api-stability}.

\begin{checklist}
  \item Betriebsannahmen dokumentieren (Rollen, Pooling, Transaction-Pooling).
  \item Context-Lebensdauern dokumentieren: Factories, Hintergrundjobs, offene Verbindungen.
  \item Mandantenübergreifender Zugriff nur über eigene Rollen, nie über App-Schalter.
  \item Bedrohungsmodell dokumentieren: was geschützt ist und was nicht
    (Seite~\pageref{sec:threat-model}).
  \item Beziehungen mandantensicher: zusammengesetzte Fremdschlüssel, eindeutig pro Mandant.
  \item Bestehende Tabellen in zwei Migrationen mit Wächter nachrüsten.
  \item Heiße Pfade (Interceptors) messen, nicht vermuten.
  \item Second-Level-Caches: der Mandant in jedem Cache-Schlüssel, ohne Mandant kein Cache.
  \item Public-API-Analyzer als Fehler, nicht als Warnung.
\end{checklist}

\section{Tiefe Erweiterungen (Teil~\ref{part:deep-extensions})}

\noindent Details: Kapitel~\ref{ch:custom-sql-expression-nodes}, \ref{ch:aggregate-translators}, \ref{ch:property-level-fluent-api-and}, \ref{ch:scaffolding-in-practice}.

\begin{checklist}
  \item Eigene SQL-Knoten: VisitChildren, Update, Quote, Print, Equals/GetHashCode.
  \item Eigene Knoten im SQL-Generator UND im Nullability-Prozessor behandeln.
  \item Aggregate: Prädikat, DISTINCT und leere Mengen bewusst behandeln.
  \item Laufzeit-Annotations bewusst wählen, Design-Time-Annotations filtern.
  \item Scaffolding: Schema-Objekte des Plugins aus der Datenbank lesen.
\end{checklist}

\section{Mehrere Provider (Teil~\ref{part:providers})}

\noindent Details: Kapitel~\ref{ch:what-belongs-in-the-core}, \ref{ch:a-second-provider-sql-server}, \ref{ch:without-sql-what-carries-over}.

\begin{checklist}
  \item Provider-neutrale Bausteine in den Kern, Unterschiede ins Provider-Paket.
  \item Relationale Bausteine in eine eigene Schicht, damit der Kern auch ohne SQL nutzbar
    bleibt.
  \item Ohne Row Level Security: dokumentieren, dass der Query-Filter die einzige Isolation ist
    und wie roher Zugriff ihn umgeht.
\end{checklist}
